\documentclass[10pt,a4paper]{amsart}
\usepackage{amsmath,amssymb,amsthm,graphicx}
\usepackage{verbatim}
\newenvironment{reference}{\comment}{\endcomment}
%\newenvironment{reference}{}{}
\newenvironment{slogan}{\comment}{\endcomment}
\newenvironment{history}{\comment}{\endcomment}
% For commutative diagrams we use Xy-pic
\usepackage[all]{xy}

% We use 2cell for 2-commutative diagrams.
\xyoption{2cell}
% Theorem environments.
%
\theoremstyle{plain}
\newtheorem{theorem}[subsection]{Theorem}
\newtheorem{proposition}[subsection]{Proposition}
\newtheorem{lemma}[subsection]{Lemma}

\theoremstyle{definition}
\newtheorem{definition}[subsection]{Definition}
\newtheorem{example}[subsection]{Example}
\newtheorem{exercise}[subsection]{Exercise}
\newtheorem{situation}[subsection]{Situation}

\theoremstyle{remark}
\newtheorem{remark}[subsection]{Remark}
\newtheorem{remarks}[subsection]{Remarks}

\numberwithin{equation}{subsection}

% Macros
%
\def\lim{\mathop{\mathrm{lim}}\nolimits}
\def\colim{\mathop{\mathrm{colim}}\nolimits}
\def\Spec{\mathop{\mathrm{Spec}}}
\def\Hom{\mathop{\mathrm{Hom}}\nolimits}
\def\Ext{\mathop{\mathrm{Ext}}\nolimits}
\def\SheafHom{\mathop{\mathcal{H}\!\mathit{om}}\nolimits}
\def\SheafExt{\mathop{\mathcal{E}\!\mathit{xt}}\nolimits}
\def\Sch{\mathit{Sch}}
\def\Mor{\mathop{\mathrm{Mor}}\nolimits}
\def\Ob{\mathop{\mathrm{Ob}}\nolimits}
\def\Sh{\mathop{\mathit{Sh}}\nolimits}
\def\NL{\mathop{N\!L}\nolimits}
\def\CH{\mathop{\mathrm{CH}}\nolimits}
\def\proetale{{pro\text{-}\acute{e}tale}}
\def\etale{{\acute{e}tale}}
\def\QCoh{\mathit{QCoh}}
\def\Ker{\mathop{\mathrm{Ker}}}
\def\Im{\mathop{\mathrm{Im}}}
\def\Coker{\mathop{\mathrm{Coker}}}
\def\Coim{\mathop{\mathrm{Coim}}}

% Boxtimes
%
\DeclareMathSymbol{\boxtimes}{\mathbin}{AMSa}{"02}

%
% Macros for moduli stacks/spaces
%
\def\QCohstack{\mathcal{QC}\!\mathit{oh}}
\def\Cohstack{\mathcal{C}\!\mathit{oh}}
\def\Spacesstack{\mathcal{S}\!\mathit{paces}}
\def\Quotfunctor{\mathrm{Quot}}
\def\Hilbfunctor{\mathrm{Hilb}}
\def\Curvesstack{\mathcal{C}\!\mathit{urves}}
\def\Polarizedstack{\mathcal{P}\!\mathit{olarized}}
\def\Complexesstack{\mathcal{C}\!\mathit{omplexes}}
% \Pic is the operator that assigns to X its picard group, usage \Pic(X)
% \Picardstack_{X/B} denotes the Picard stack of X over B
% \Picardfunctor_{X/B} denotes the Picard functor of X over B
\def\Pic{\mathop{\mathrm{Pic}}\nolimits}
\def\Picardstack{\mathcal{P}\!\mathit{ic}}
\def\Picardfunctor{\mathrm{Pic}}
\def\Deformationcategory{\mathcal{D}\!\mathit{ef}}
\usepackage{cleveref}
\crefname{originallemma}{original lemma}{original lemmas}
\Crefname{originallemma}{Original lemma}{Original lemmas}
\makeatletter
\newlabel{lemma-hom-finite-presentation-colimit}{{lemma-hom-finite-presentation-colimit}{source}{}{sourcepointer-1}{}}
\newlabel{lemma-hom-finite-presentation-colimit@cref}{{[originallemma][0][]lemma-hom-finite-presentation-colimit}{[1][0][]source}{}{}{}}
\makeatother
\crefname{originalreference}{original reference}{original references}
\Crefname{originalreference}{Original reference}{Original references}
\makeatletter
\newlabel{categories-section-posets-limits}{{categories-section-posets-limits}{source}{}{sourcepointer-2}{}}
\newlabel{categories-section-posets-limits@cref}{{[originalreference][0][]categories-section-posets-limits}{[1][0][]source}{}{}{}}
\makeatother
\crefname{originalreference}{original reference}{original references}
\Crefname{originalreference}{Original reference}{Original references}
\makeatletter
\newlabel{sheaves-lemma-directed-colimits-sections}{{sheaves-lemma-directed-colimits-sections}{source}{}{sourcepointer-3}{}}
\newlabel{sheaves-lemma-directed-colimits-sections@cref}{{[originalreference][0][]sheaves-lemma-directed-colimits-sections}{[1][0][]source}{}{}{}}
\makeatother
\title{Homomorphisms out of a finitely presented module commute with filtered colimits}
\author{The Stacks Project Authors}
\date{stacks-project:modules.tex:local-snapshot-20261012; GNU FDL 1.2 or later}
\begin{document}
\maketitle
\noindent\textit{Editorial note.} All delivered mathematical text is the author's own, copied from the pinned Stacks Project chapter source sources/stacks-modules/pinned-original-source.tex; no statement or proof step was rewritten, reordered or removed.\par\medskip
\noindent\textit{Editorial note.} Preamble provenance: the Stacks Project keeps one single shared preamble for every chapter; this card takes it directly from the staged sources/stacks-modules/preamble.tex as three verbatim spans (lines 8-12, 14-18 and 154-217), so no derived file was written and nothing was added or redefined. The wrapper class and the base packages are supplied by the assembler.\par\medskip
\noindent\textit{Editorial note.} Bibliography provenance: neither the counted statement nor its proof contains a citation (this lane rescans the whole span before writing the card), so the card carries no bibliography and no entry was reconstructed.\par\medskip
\noindent\textit{Editorial note.} Selection provenance: this card was opened under the file-based control contract (CONTROL.md), with the mechanical pre-card gate run on its statement and proof interval before the card was written.\par\medskip
\noindent\textit{Allowed standard prerequisites (editorial).}\begin{itemize}
\item The description of morphisms of modules as the global sections of the sheaf of homomorphisms. Recalled inside the counted proof (source lines 4062-4067). Turns the two sides of the displayed equality into global sections of sheaves of homomorphisms.
\item The commutation of the sheaf of homomorphisms out of a finitely presented module with filtered colimits, and the compatibility of global sections with directed colimits under a finiteness hypothesis on the coverings. Both are invoked at the end of the counted proof (source lines 4068-4070) as a pointer to a lemma of this chapter and a cross-chapter pointer. Identify the colimit of the sheaves of homomorphisms with the sheaf of homomorphisms into the colimit and then take global sections.
\end{itemize}\medskip
\ifdefined\hypertarget\hypertarget{sourcepointer-1}{}\fi
\noindent\textit{Source reference.} Original lemma lemma-hom-finite-presentation-colimit: sources/stacks-modules/pinned-original-source.tex, source lines 4068--4068. Navigation pointer only: the counted text refers to the Stacks item lemma-hom-finite-presentation-colimit, which belongs to another chapter of the project and is not part of this card. The pointer carries the target's own label as its number because this chapter's staged snapshot does not contain the other chapter's tag table; its content is not reproduced here.\par\medskip
\ifdefined\hypertarget\hypertarget{sourcepointer-2}{}\fi
\noindent\textit{Source reference.} Original reference categories-section-posets-limits: sources/stacks-modules/pinned-original-source.tex, source lines 1148--1148. Navigation pointer only: the counted text refers to the Stacks item categories-section-posets-limits, which belongs to another chapter of the project and is not part of this card. The pointer carries the target's own label as its number because this chapter's staged snapshot does not contain the other chapter's tag table; its content is not reproduced here.\par\medskip
\ifdefined\hypertarget\hypertarget{sourcepointer-3}{}\fi
\noindent\textit{Source reference.} Original reference sheaves-lemma-directed-colimits-sections: sources/stacks-modules/pinned-original-source.tex, source lines 3271--3271. Navigation pointer only: the counted text refers to the Stacks item sheaves-lemma-directed-colimits-sections, which belongs to another chapter of the project and is not part of this card. The pointer carries the target's own label as its number because this chapter's staged snapshot does not contain the other chapter's tag table; its content is not reproduced here.\par\medskip
\begin{lemma}
\label{lemma-finite-presentation-quasi-compact-colimit}
Let $X$ be a ringed space.
Let $I$ be a preordered set and
let $(\mathcal{F}_i, \varphi_{ii'})$ be a system over $I$
consisting of sheaves of $\mathcal{O}_X$-modules
(see Categories, Section \ref{categories-section-posets-limits}).
Assume
\begin{enumerate}
\item $I$ is directed,
\item $\mathcal{G}$ is an $\mathcal{O}_X$-module of finite presentation, and
\item $X$ has a cofinal system of open coverings
$\mathcal{U} : X = \bigcup_{j\in J} U_j$ with
$J$ finite and $U_j \cap U_{j'}$ quasi-compact
for all $j, j' \in J$.
\end{enumerate}
Then we have
$$
\colim_i \Hom_X(\mathcal{G}, \mathcal{F}_i)
=
\Hom_X(\mathcal{G}, \colim_i \mathcal{F}_i).
$$
\end{lemma}

\begin{proof}
Set
$\mathcal{H} = \SheafHom_{\mathcal{O}_X}(\mathcal{G}, \colim \mathcal{F}_i)$
and $\mathcal{H}_i = \SheafHom_{\mathcal{O}_X}(\mathcal{G}, \mathcal{F}_i)$.
Recall that
$$
\Hom_X(\mathcal{G}, \mathcal{F}) = \Gamma(X, \mathcal{H})
\quad\text{and}\quad
\Hom_X(\mathcal{G}, \mathcal{F}_i) = \Gamma(X, \mathcal{H}_i)
$$
by construction. By Lemma \ref{lemma-hom-finite-presentation-colimit} we have
$\mathcal{H} = \colim \mathcal{H}_i$. Thus the lemma follows from
Sheaves, Lemma \ref{sheaves-lemma-directed-colimits-sections}.
\end{proof}

\end{document}
